\section{Introduction}
\label{sec:introduction}

Bounded decisions are classifications with a declared label set: an intent, a toxicity bit, a spam bit, a clause type, a yes/no/maybe answer, a fact-checking label, or a product-relevance label. The same decision can be implemented in more than one way. A rule matches a pattern. A classifier is fit on labeled text. A typed decision model returns a distribution over a label set supplied at inference time. A generative model writes text that a parser then accepts or rejects as a label.

Algorithm selection already treats the choice of procedure as a map from problem features to performance~\citep{rice1976,smithmiles2009}. LLM routing already chooses among language models, often with cost in view, and sometimes with a notion of task type (Shnitzer et al., \citeyear{shnitzer2023}; Ong et al., \citeyear{ong2024}; Liu et al., \citeyear{liu2026trouter}; Li, Zhang, et al., \citeyear{li2026llmrouterbench}). Comparisons of encoders or linear classifiers with prompted generative models already report quality together with latency or dollars on fixed-label text tasks~\citep{zhang2025tamas,gonzalez2026,bucher2024,vajjala2025}. Leave-one-dataset-out selection from training-split meta-features has been studied for abstractive summarization~\citep{bino2026}. This pilot does not claim those problems as new.

What we measure is narrower. Each task is scored, under one input string and one label set, by the primitives that the protocol marks as eligible. A fingerprint is computed from the task contract and the training split only. A ridge model, specified before these predictive results, maps that fingerprint and a primitive identity to the measured profile. We then ask whether the predicted profile is closer to the measured profile than the training-fold mean of that primitive, and whether any closer prediction changes the primitive selected by a policy that was also specified in advance.

The independent unit for that second question is the task. The evaluation rows inside a dataset estimate that dataset's scores. They do not multiply the number of tasks. With eight tasks, and with three of the five selector-training families containing a single task, the pilot cannot support a general claim about family-level transfer. We state that limit in the body of the paper.

The empirical reading is as follows. Computational primitives exhibit substantial task-dependent differences in quality and operational behavior, but in this eight-task pilot, a simple pre-deployment task fingerprint did not reliably improve prediction over primitive-specific training-fold mean baselines or alter the resulting computational choice. This does not say that the hypothesis is false, that the classifier is universally best, or that we have introduced a router.

\section{Research questions}
\label{sec:rqs}

\begin{enumerate}[label=RQ\arabic*.]
\item How heterogeneous are quality, validity, latency, and cost across computational primitives on these bounded decision tasks?
\item How much predictive information about primitive behavior, in this pilot, is contained in the fixed pre-deployment fingerprint?
\item Does that fingerprint, through the pre-specified ridge model, improve held-out prediction over the primitive-specific mean?
\item Does a closer predicted profile change the primitive selected by the pre-specified policies?
\item What failure modes appear when the pinned generative configuration is used as a bounded decision primitive?
\end{enumerate}

\noindent RQ2 and RQ3 are about this fingerprint and this model. They are not questions about every possible fingerprint. RQ1 and RQ5 are descriptive. RQ4 separates prediction error from the downstream choice.
